\chapter{HAI Robotics K50 Reference}

\section*{AMR Inspired by HAI Robotics K50 Fast-Transit Companion System :}

\begin{figure}[H]
\centering
\fbox{\parbox{0.85\textwidth}{
\centering
\textbf{Image placeholder}\\
Insert the rack interaction reference image here.
}}
\caption{Rack interaction reference image placeholder.}
\end{figure}

\section*{Quick and convenient AMRs feature:}

\begin{itemize}
    \item A max speed of 4m/s
\end{itemize}

\begin{itemize}
    \item Maximum jacking height of 0.75m suitable for ergonomic workstations
\end{itemize}

\begin{itemize}
    \item LiDAR remote scanning
\end{itemize}

\section*{Mechanical specifications}

Dimensions (L $\times$ W $\times$ H) (mm) : 658$\times$460$\times$330 Net weight (kg) : 81 Rated load (kg) : 30/50 Case type mode : Carton, tote, and other materials with regular shapes Driving mode : Differential-wheeled Navigation DM code + IMU

\section*{Motion performance}

Driving speed (m/s) $\leq$ 4.0 Stop accuracy (mm) $\pm$ 10 Travel direction: Forward and backward + rotation in place + turning

\section*{Functional Description}

The proposed AMR system is inspired by K50 Fast-Transit Companion AMR developed by HAI Robotics. The robot is designed to operate as a fast-transit mobile platform capable of interacting directly with standardized warehouse racks and handling totes through a structured ``backpack'' mechanism.

The core operational concept is based on a three-step rack interaction workflow, enabling efficient retrieval, transport, and delivery of storage totes.

a) Precise Positioning Under Rack

The AMR navigates autonomously within warehouse aisles and performs precise alignment beneath storage racks equipped with low-height guide rails (approximately 0.43 m from the ground).

Once positioned correctly, the robot aligns itself under the target rack bay where totes are pre-stored by an upper system or external automation layer.

\section*{This phase requires:}

\begin{itemize}
    \item High-precision localization
    \item Fine motion control near racks
    \item Accurate rack alignment detection (visual/LiDAR-based)
\end{itemize}

b) Tote Securing (Backpack Mode Mechanism)

After correct positioning, the AMR engages a mechanical lifting and securing mechanism composed of twin cradling supports or fork-like arms.

\begin{figure}[H]
\centering
\fbox{\parbox{0.85\textwidth}{
\centering
\textbf{Image placeholder}\\
Insert the tote securing mechanism reference image here.
}}
\caption{Tote securing mechanism reference image placeholder.}
\end{figure}

\section*{This mechanism:}

\begin{itemize}
    \item Engages the tote from the underside
    \item Stabilizes the load laterally
    \item Locks the tote securely onto the robot chassis
    \item Transforms the tote into a ``backpack-style carried load''
\end{itemize}

At this stage, the tote becomes mechanically integrated with the robot, ensuring stability during high-speed transport.

c) High-Speed Transport and Delivery to Workstation

Once the tote is secured, the AMR transitions into fast-transit mode, transporting the loaded tote to a designated workstation or drop-off zone.

\section*{Upon arrival, the robot:}

\begin{itemize}
    \item Navigates to a predefined delivery point
    \item Raises or positions the tote at an ergonomically optimized height for human
operators

    \item Allows efficient manual picking or unloading operations
\end{itemize}

This ensures smooth human-robot collaboration by minimizing worker bending and improving picking efficiency.

\begin{figure}[H]
\centering
\fbox{\parbox{0.85\textwidth}{
\centering
\textbf{Image placeholder}\\
Insert the delivery workflow reference image here.
}}
\caption{Delivery workflow reference image placeholder.}
\end{figure}

\section*{System Behavior Summary}

\section*{The Fast-Transit Companion AMR operates as a rack-interacting logistics unit that:}

\begin{itemize}
    \item Locates and aligns under warehouse racks
    \item Engages and secures totes mechanically (``backpack mode'')
    \item Transports loads efficiently across warehouse zones
    \item Delivers totes at ergonomic workstation heights for operators
\end{itemize}

These robots operate in coordination with ACR (Automated Case-handling Robots) that utilize a grappling-hook-based lifting mechanism. This design enables the ACR systems to directly engage and retrieve totes from storage racks configured in single-, double-, or triple-deep layouts.

The AMR and ACR systems work as a complementary logistics pair: the ACR handles high-
